\Needspace{9\baselineskip}\begin{center}\textit{Eine Methode, den Grad der Gleichungen für trigonometrische Funktionen schrittweise zu senken.}\end{center}


\Needspace{6\baselineskip}\Article{363}
Wie sich unter der Voraussetzung \(n-1=ef\) die Funktion \(X\) in \(e\) Faktoren vom Grad \(f\) zerlegen lässt, sobald die Werte aller \(e\) Summen von \(f\) Gliedern bekannt sind (Artikel 348), so lässt sich dann auch, wenn man voraussetzt, dass \(Z=0\) eine Gleichung vom Grad \(n-1\) ist, deren Wurzeln die Sinus oder beliebige andere trigonometrische Funktionen der Winkel \(\frac Pn,\allowbreak \frac{2P}n\ldots\frac{(n-1)P}n\) sind, die Funktion \(Z\) in \(e\) Faktoren vom Grad \(f\) zerlegen. Die Hauptpunkte dieses Verfahrens sind die folgenden.

Es bestehe \(\Omega\) aus den \(e\) Perioden von \(f\) Gliedern \((f,1)=P,\allowbreak P'.P''\), usw.; die Periode \(P\) aus den Wurzeln \([1],\allowbreak [a],\allowbreak [b],\allowbreak [c]\), usw.; \(P'\) aus \([a'],\allowbreak [b'],\allowbreak [c']\), usw.; \(P''\) aus \([a''],\allowbreak [b''],\allowbreak [c'']\), usw., usw. Der Wurzel \([1]\) entspreche der Winkel \(\omega\), also den Wurzeln \([a],\allowbreak [b]\), usw., die Winkel \(a\omega,\allowbreak b\omega\), usw., den Wurzeln \([a'],\allowbreak [b']\), usw., die Winkel \(a'\omega,\allowbreak b'\omega\), usw., und den Wurzeln \([a''],\allowbreak [b'']\), usw., die Winkel \(a''\omega,\allowbreak b''\omega\), usw., usw. Dann erkennt man leicht, dass alle diese Winkel zusammengenommen mit den Winkeln \(\frac Pn,\allowbreak \frac{2P}n,\allowbreak \frac{3P}n\ldots\frac{(n-1)P}n\) hinsichtlich der trigonometrischen Funktionen\HistoricalNote{*)}{In dieser Hinsicht stimmen zwei Winkel überein, deren Differenz dem ganzen Kreisumfang oder irgendeinem Vielfachen davon gleich ist. Solche Winkel könnten wir \textit{nach dem Kreisumfang kongruent} nennen, wenn man die Kongruenz in einem etwas weiteren Sinne verstehen wollte.} übereinstimmen. Bezeichnet man daher die betreffende Funktion durch das dem Winkel vorangestellte Zeichen \(\varphi\) und setzt das Produkt aus den \(e\) Faktoren \[x-\varphi\omega,\quad x-\varphi a\omega,\quad x-\varphi b\omega\quad\text{usw.}\] \(=Y\), das Produkt aus \(x-\varphi a'\omega,\allowbreak x-\varphi b'\omega\), usw., \(=Y'\), das Produkt aus \(x-\varphi a''\omega,\allowbreak x-\varphi b''\omega\), usw., \(=Y''\), und so weiter, so ist notwendig das Produkt \(YY'Y''\ldots=Z\). Nun bleibt zu zeigen, dass sich alle Koeffizienten in den Funktionen \(Y,\allowbreak Y',\allowbreak Y''\), usw., auf eine Form wie
\SourceLeaf{463}
\[A+B(f,1)+C(f,g)+D(f,gg)\ldots+L(f,g^{e-1})\] zurückführen lassen. Danach sind offenbar alle als bekannt anzusehen, sobald die Werte aller Summen von \(f\) Gliedern bekannt sind. Dies werden wir folgendermaßen erreichen.

Wie \(\cos\omega=\tfrac12[1]+\tfrac12[1]^{n-1},\allowbreak \ \sin\omega=-\tfrac12i[1]+\tfrac12i[1]^{n-1}\), so lassen sich nach dem vorigen Artikel auch die übrigen trigonometrischen Funktionen des Winkels \(\omega\) auf eine Form wie \(\mathfrak A+\mathfrak B[1]+\mathfrak C[1]^2+\mathfrak D[1]^3+\text{usw.}\) zurückführen. Man erkennt ohne Schwierigkeit, dass die Funktion des Winkels \(k\omega\) dann \(=\mathfrak A+\mathfrak B[k]+\mathfrak C[k]^2+\mathfrak D[k]^3+\text{usw.}\) wird, wobei \(k\) eine beliebige ganze Zahl bezeichnet. Da nun jeder Koeffizient in \(Y\) ein symmetrisches Polynom in \(\varphi\omega,\allowbreak \varphi a\omega,\allowbreak \varphi b\omega\), usw., ist, ist klar: Setzt man für diese Größen ihre Werte ein, so wird jeder Koeffizient ein symmetrisches Polynom in \([1],\allowbreak [a],\allowbreak [b]\), usw.; daher lassen sie sich nach Artikel 347 auf die Form \(A+B(f,1)+C(f,g)+\text{usw.}\) zurückführen. Mit genau derselben Begründung kann man auch alle Koeffizienten in \(Y',\allowbreak Y''\), usw., auf eine ähnliche Form zurückführen. \textit{Q. E. D.}

\Needspace{6\baselineskip}\Article{364}
Zum Problem des vorigen Artikels fügen wir noch einige Beobachtungen hinzu.

I. Da die einzelnen Koeffizienten in \(Y'\) Funktionen der in der Periode \(P'\) (die man \(=(f,a')\) setzen darf) enthaltenen Wurzeln sind, die mit den durch die entsprechenden Koeffizienten gegebenen Funktionen der Wurzeln in \(P\) in \(Y\) übereinstimmen, ist aus Artikel 347 klar, dass sich \(Y'\) aus \(Y\) ableiten lässt, wenn man überall in \(Y\) für \((f,1),\allowbreak (f,g),\allowbreak (f,gg)\), usw., der Reihe nach \((f,a'),\allowbreak (f,a'g),\allowbreak (f,a'gg)\), usw., einsetzt. Ebenso wird \(Y''\) aus \(Y\) abgeleitet, indem man überall in \(Y\) für \((f,1),\allowbreak (f,g),\allowbreak (f,gg)\), usw., der Reihe nach \((f,a''),\allowbreak (f,a''g),\allowbreak (f,a''gg)\), usw., usw., einsetzt. Sobald also die Funktion \(Y\) entwickelt ist, folgen die übrigen Funktionen \(Y',\allowbreak Y''\), usw., daraus ohne Schwierigkeit.

II. Setzt man \[Y=x^f-\alpha x^{f-1}+\beta x^{f-2}-\text{usw.}\] voraus, so sind die Koeffizienten \(\alpha,\allowbreak \beta\), usw., der Reihe nach die Summe der Wurzeln der Gleichung \(Y=0\quad \text{d.\,h.}\) der Größen \(\varphi\omega,\allowbreak \varphi a\omega,\allowbreak \varphi b\omega\), usw., die Summe der Produkte aus je zweien und so weiter. Meist erhält man diese Koeffizienten aber viel bequemer mit einer der in Artikel 349 angegebenen ähnlichen Methode, indem man die Summe der Wurzeln \(\varphi\omega,\allowbreak \varphi a\omega,\allowbreak \varphi b\omega\), usw., die Summe ihrer Quadrate,
\SourceLeaf{464}
Kuben, usw., berechnet und daraus nach dem Newtonschen Satz jene Koeffizienten ableitet. — Bezeichnet \(\varphi\) den Tangens, Sekans, Kotangens oder Kosekans, so gibt es noch weitere Abkürzungen, die wir hier jedoch ohne Erörterung übergehen.

III. Besondere Betrachtung verdient der Fall, dass \(f\) gerade ist und daher jede Periode \(P,\allowbreak P',\allowbreak P''\), usw., aus \(\tfrac12f\) Perioden von je zwei Gliedern besteht. Es bestehe \(P\) aus \((2,1),\allowbreak (2,\mathfrak{a}),\allowbreak (2,\mathfrak{b}),\allowbreak (2,\mathfrak{c})\), usw. Dann stimmen die Zahlen \(1,\allowbreak \mathfrak{a},\allowbreak \mathfrak{b},\allowbreak \mathfrak{c}\), usw., und \(n-1,\allowbreak n-\mathfrak{a}.n-\mathfrak{b},\allowbreak n-\mathfrak{c}\), usw., zusammengenommen mit \(1,\allowbreak a,\allowbreak b,\allowbreak c\), usw., überein oder sind ihnen wenigstens (was hier auf dasselbe hinausläuft) modulo \(n\) kongruent. Es gilt aber \(\varphi(n-1)\omega=\pm\varphi\omega,\allowbreak \ \varphi(n-\mathfrak{a})\omega=\pm\varphi \mathfrak{a}\omega\), usw., wobei die oberen Vorzeichen gelten, wenn \(\varphi\) den Kosinus oder Sekans bezeichnet, die unteren, wenn \(\varphi\) den Sinus, Tangens, Kotangens oder Kosekans bezeichnet. Daraus folgt, dass in den ersten beiden Fällen unter den Faktoren, aus denen \(Y\) besteht, stets je zwei gleich sind und daher \(Y\) ein Quadrat ist, und zwar \(Y=yy\), wenn \(y\) gleich dem folgenden Produkt gesetzt wird: \[x-\varphi\omega,\ x-\varphi \mathfrak{a}\omega,\ x-\varphi \mathfrak{b}\omega\text{ usw.}\]

Ebenso sind in denselben Fällen die übrigen Funktionen \(Y',\allowbreak Y''\), usw., Quadrate. Setzt man nämlich voraus, dass \(P'\) aus \((2,\mathfrak{a}'),\allowbreak (2,\mathfrak{b}'),\allowbreak (2,\mathfrak{c}')\), usw., \(P''\) aus \((2,\mathfrak{a}''),\allowbreak (2,\mathfrak{b}''),\allowbreak (2,\mathfrak{c}'')\), usw., usw., besteht und dass das Produkt aus \(x-\varphi \mathfrak{a}'\omega,\allowbreak x-\varphi \mathfrak{b}'\omega,\allowbreak x-\varphi \mathfrak{c}'\omega\), usw., \(=y'\), das Produkt aus \(x-\varphi \mathfrak{a}''\omega,\allowbreak x-\varphi \mathfrak{b}''\omega\), usw., \(=y''\) ist, und so weiter, so wird \(Y'=y'y',\allowbreak \ Y''=y''y''\), usw. Auch die Funktion \(Z\) ist dann ein Quadrat (vgl. oben Artikel 337), dessen Quadratwurzel gleich dem Produkt aus \(y,\allowbreak y',\allowbreak y''\), usw., ist. Ferner erkennt man leicht, dass \(y',\allowbreak y''\), usw., ebenso aus \(y\) abgeleitet werden, wie wir oben in I gesagt haben, dass \(Y',\allowbreak Y''\), usw., aus \(Y\) folgen. Auch jeder Koeffizient in \(y\) lässt sich auf die Form \[A+B(f,1)+C(f,g)+\text{usw.}\] zurückführen, da die Summen der einzelnen Potenzen der Wurzeln der Gleichung \(y=0\) offenbar die Hälften der Potenzen der Wurzeln der Gleichung \(Y=0\) sind und sich somit auf eine solche Form zurückführen lassen. In den letzten vier Fällen dagegen ist \(Y\) das Produkt aus den Faktoren \[xx-(\varphi\omega)^2,\ xx-(\varphi \mathfrak{a}\omega)^2,\ xx-(\varphi \mathfrak{b}\omega)^2\text{ usw.}\] und hat daher die Form \[x^f-\lambda x^{f-2}+\mu x^{f-4}-\text{usw.}\] Offenbar lassen sich die Koeffizienten \(\lambda,\allowbreak \mu\), usw., aus den Summen der Quadrate, vierten Potenzen, usw., der Wur\SourceLeaf{465}zeln \(\varphi\omega,\allowbreak \varphi \mathfrak{a}\omega,\allowbreak \varphi \mathfrak{b}\omega\), usw., ableiten; und ebenso verhalten sich die Funktionen \(Y',\allowbreak Y''\), usw.

\textit{Beispiel} I. Es sei \(n=17,\allowbreak f=8\), und \(\varphi\) bezeichne den Kosinus. Daraus folgt \[\begin{aligned}Z=\bigl(&x^8+\tfrac12x^7-\tfrac74x^6-\tfrac34x^5+\tfrac{15}{16}x^4\\&+\tfrac5{16}x^3-\tfrac5{32}xx-\tfrac1{32}x+\tfrac1{256}\bigr)^2\end{aligned}\] und man muss daher \(\surd Z\) in zwei Faktoren vierten Grades \(y,\allowbreak y'\) zerlegen. Die Periode \(P=(8,1)\) besteht aus \((2,1),\allowbreak (2,9),\allowbreak (2,13),\allowbreak (2,15)\), weshalb \(y\) das Produkt aus den Faktoren \[x-\varphi\omega,\ x-\varphi9\omega,\ x-\varphi13\omega,\ x-\varphi15\omega\] ist. Setzt man \(\tfrac12[k]+\tfrac12[n-k]\) für \(\varphi k\omega\) ein, so findet man \[\begin{gathered}\varphi\omega+\varphi9\omega+\varphi13\omega+\varphi15\omega=\tfrac12(8,1),\\(\varphi\omega)^2+(\varphi9\omega)^2+(\varphi13\omega)^2+(\varphi15\omega)^2=2+\tfrac14(8,1)\end{gathered}\] Ebenso ist die Summe der Kuben \(=\tfrac38(8,1)+\tfrac18(8,3)\), die Summe der vierten Potenzen \(=1\tfrac12+\tfrac5{16}(8,1)\). Bestimmt man hieraus nach dem Newtonschen Satz die Koeffizienten in \(y\), so ergibt sich \[\begin{aligned}y&=x^4-\tfrac12(8,1)x^3+\tfrac14((8,1)+2(8,3))xx\\&\quad-\tfrac18((8,1)+3(8,3))x+\tfrac1{16}((8,1)+(8,3))\end{aligned}\] Dagegen wird \(y'\) aus \(y\) abgeleitet, indem man \((8,1)\) mit \((8,3)\) vertauscht. Setzt man also für \((8,1),\allowbreak (8,3)\) die Werte \(-\tfrac12+\tfrac12\surd17,\allowbreak -\tfrac12-\tfrac12\surd17\) ein, so ergibt sich \[\begin{aligned}y&=x^4+(\tfrac14-\tfrac14\surd17)x^3-(\tfrac38+\tfrac18\surd17)xx\\&\quad+(\tfrac14+\tfrac18\surd17)x-\tfrac1{16}\end{aligned}\] \[\begin{aligned}y'&=x^4+(\tfrac14+\tfrac14\surd17)x^3-(\tfrac38-\tfrac18\surd17)xx\\&\quad+(\tfrac14-\tfrac18\surd17)x-\tfrac1{16}\end{aligned}\] Ebenso lässt sich \(\surd Z\) in vier Faktoren zweiten Grades zerlegen, deren erster \((x-\varphi\omega)(x-\varphi13\omega)\), zweiter \((x-\varphi9\omega)(x-\varphi15\omega)\), dritter \((x-\varphi3\omega)(x-\varphi5\omega)\) und vierter \((x-\varphi10\omega)(x-\varphi11\omega)\) ist. Alle Koeffizienten dieser Faktoren können durch die vier Summen \((4,1),\allowbreak (4,9),\allowbreak (4,3),\allowbreak (4,10)\) ausgedrückt werden. Offenbar ist aber das Produkt aus dem ersten und zweiten Faktor \(y\), das Produkt aus dem dritten und vierten \(y'\).

\textit{Beispiel} II. Soll bei unveränderten übrigen Voraussetzungen \(\varphi\) den Sinus bezeichnen, sodass \[\begin{aligned}Z&=x^{16}-\tfrac{17}4x^{14}+\tfrac{119}{16}x^{12}-\tfrac{221}{32}x^{10}\\&\quad+\tfrac{935}{256}x^8-\tfrac{561}{512}x^6+\tfrac{357}{2048}x^4-\tfrac{51}{4096}xx+\tfrac{17}{65536}\end{aligned}\] in zwei Faktoren achten Grades \(y,\allowbreak y'\) zu zerlegen ist, so ist \(y\) das Produkt aus
\SourceLeaf{466}
vier quadratischen Faktoren \[xx-(\varphi\omega)^2,\ xx-(\varphi9\omega)^2,\ xx-(\varphi13\omega)^2,\ xx-(\varphi15\omega)^2\] Da nun \(\varphi k\omega=-\tfrac12i[k]+\tfrac12i[n-k]\), wird \[\begin{aligned}(\varphi k\omega)^2&=-\tfrac14[2k]+\tfrac12[n]-\tfrac14[2n-2k]\\&=\tfrac12-\tfrac14[2k]-\tfrac14[2n-2k]\end{aligned}\] Hieraus ergibt sich für die Summe der Quadrate der Wurzeln \(\varphi\omega,\allowbreak \varphi9\omega,\allowbreak \varphi13\omega.\varphi15\omega\) der Wert \(2-\tfrac14(8,1)\), für die Summe ihrer vierten Potenzen \(=\tfrac32-\tfrac3{16}(8,1)\), für die Summe ihrer sechsten Potenzen \(=\tfrac54-\tfrac9{64}(8,1)-\tfrac1{64}(8,3)\) und für die Summe ihrer achten Potenzen \(\tfrac{35}{32}-\tfrac{27}{256}(8,1)-\tfrac1{32}(8,3)\). Daraus folgt \[\begin{aligned}y={}&x^8-(2-\tfrac14(8,1))x^6\\&+(\tfrac32-\tfrac5{16}(8,1)+\tfrac18(8,3))x^4\\&-(\tfrac12-\tfrac9{64}(8,1)+\tfrac5{64}(8,3))xx\\&+\tfrac1{16}-\tfrac5{256}(8,1)+\tfrac3{256}(8,3)\end{aligned}\] Man leitet \(y'\) aus \(y\) durch Vertauschung von \((8,1),\allowbreak (8,3)\) ab, sodass man durch Einsetzen der Werte dieser Summen \[\begin{aligned}y&=x^8-(\tfrac{17}8-\tfrac18\surd17)x^6+(\tfrac{51}{32}-\tfrac7{32}\surd17)x^4\\&\quad-(\tfrac{17}{32}-\tfrac7{64}\surd17)xx+\tfrac{17}{256}-\tfrac1{64}\surd17\end{aligned}\] \[\begin{aligned}y'&=x^8-(\tfrac{17}8+\tfrac18\surd17)x^6+(\tfrac{51}{32}+\tfrac7{32}\surd17)x^4\\&\quad-(\tfrac{17}{32}+\tfrac7{64}\surd17)xx+\tfrac{17}{256}+\tfrac1{64}\surd17\end{aligned}\] erhält. Ebenso lässt sich \(Z\) in vier Faktoren zerlegen, deren Koeffizienten durch Summen von vier Gliedern ausgedrückt werden können; dabei ist das Produkt aus zweien \(y\) und das Produkt aus den beiden übrigen \(y'\).


